\documentclass[conference]{IEEEtran}
\IEEEoverridecommandlockouts
% The preceding line is only needed to identify funding in the first footnote. If that is unneeded, please comment it out.
\usepackage{cite}
\usepackage{amsmath,amssymb,amsfonts}
\usepackage{algorithmic}
\usepackage{graphicx}
\usepackage{textcomp}
\usepackage{xcolor}
\def\BibTeX{{\rm B\kern-.05em{\sc i\kern-.025em b}\kern-.08em
    T\kern-.1667em\lower.7ex\hbox{E}\kern-.125emX}}

\begin{document}

\title{Bidirectional Range-Azimuth Upsampling for 1-Bit SAR Imaging
\thanks{Identify applicable funding agency here. If none, delete this.}
}

\author{\IEEEauthorblockN{1\textsuperscript{st} Given Name Surname}
\IEEEauthorblockA{\textit{dept. name of organization (of Aff.)} \\
\textit{name of organization (of Aff.)}\\
City, Country \\
email address or ORCID}
~\\
\and
\IEEEauthorblockN{2\textsuperscript{nd} Given Name Surname*}
\IEEEauthorblockA{\textit{dept. name of organization (of Aff.)} \\
\textit{name of organization (of Aff.)}\\
City, Country \\
email address or ORCID}
*Corresponding author
~\\
\and
\IEEEauthorblockN{3\textsuperscript{rd} Given Name Surname}
\IEEEauthorblockA{\textit{dept. name of organization (of Aff.)} \\
\textit{name of organization (of Aff.)}\\
City, Country \\
email address or ORCID}
}

\maketitle

\begin{abstract}
One-bit synthetic aperture radar (SAR) imaging reduces sampling precision by retaining only sign information, but the resulting quantization distortion can severely degrade image quality. Existing studies mainly improve one-bit SAR imaging through threshold design or post-processing. This paper studies a complementary allocation question: under a fixed total upsampling budget, should the budget be assigned to only range or azimuth, or split across both dimensions? We formulate bidirectional range-azimuth upsampling (BRAU), where the total budget is fixed as $Q=R\times A$ and the range and azimuth factors are jointly selected. Experiments on seven SAR datasets show that, for all tested integer budgets, the best bidirectional allocation consistently outperforms the best unidirectional allocation in PSNR and SSIM. The advantage also holds for non-integer factors and under zero-threshold, non-random constant-threshold, and random-threshold quantization. Spectral analysis further shows that bidirectional allocation produces lower off-support and directional leakage after range compression and range cell migration correction. These results indicate that splitting the upsampling budget across range and azimuth is a simple and effective design principle for one-bit SAR imaging.
\end{abstract}

\begin{IEEEkeywords}
one-bit SAR imaging, upsampling, range-azimuth allocation, quantization, spectral leakage
\end{IEEEkeywords}

\section{Introduction}

Synthetic aperture radar (SAR) forms high-resolution images by coherently processing echoes collected along a synthetic aperture, and has become a core sensing modality for remote sensing and target observation \cite{cumming2005sar,curlander1991sar}. Conventional SAR systems usually rely on high-resolution analog-to-digital converters to preserve echo amplitude and phase information. On lightweight or resource-constrained platforms, however, the cost, power consumption, and data throughput of high-resolution sampling can become limiting factors. One-bit SAR imaging addresses this constraint by retaining only the sign of the received echo, which substantially simplifies the sampling hardware but also removes amplitude information and introduces strong quantization distortion.

Early signum-coded SAR studies showed that meaningful SAR images can still be formed from sign-only data under suitable processing conditions \cite{franceschetti1991signum}. Recent one-bit SAR methods further improve image quality by designing quantization thresholds or reconstruction strategies. Single-frequency thresholds were shown to suppress harmonic interference in one-bit SAR imaging \cite{zhao2019onebit}, while time-varying and fixed-threshold strategies have been investigated to improve the fidelity of low-precision SAR echoes \cite{demir2018onebit,zhao2021lowprecision,nie2025fixed}. Learning-based de-quantization methods have also been introduced to suppress one-bit harmonics after data acquisition \cite{si2023dequantization}. These studies demonstrate that threshold design and post-processing can recover useful information from severely quantized SAR data.

Most existing work focuses on how to construct thresholds or how to recover images from one-bit measurements. In contrast, this paper studies a simpler but less explored design question: under a fixed total upsampling budget, should the budget be concentrated in one dimension, or split between the range and azimuth dimensions? Let $R$ and $A$ denote the range and azimuth upsampling factors, respectively, and let $Q=R\times A$ be the total upsampling budget. A unidirectional allocation uses either $R=Q,A=1$ or $R=1,A=Q$, whereas a bidirectional allocation uses $R>1$ and $A>1$ under the same $Q$.

This paper presents bidirectional range-azimuth upsampling (BRAU) as a fixed-budget allocation strategy for one-bit SAR imaging. Experiments on seven SAR datasets show that, for all tested integer budgets, the best bidirectional allocation consistently outperforms the best unidirectional allocation in both PSNR and SSIM. The advantage also appears under non-integer upsampling factors and remains significant under zero-threshold, non-random constant-threshold, and random-threshold settings. Mechanism analysis further shows that bidirectional allocation produces lower off-support and directional spectral leakage after range compression, supporting the interpretation that range-azimuth allocation better matches the two-dimensional spectral structure of SAR echoes.

The main contributions are summarized as follows:
\begin{itemize}
    \item We formulate fixed-budget range-azimuth upsampling allocation for one-bit SAR imaging, where the total budget $Q=R\times A$ is kept constant.
    \item We show that bidirectional allocation consistently improves reconstruction quality over the best unidirectional allocation across multiple tested budgets.
    \item We verify that the advantage is robust to non-integer upsampling factors and different threshold constructions.
    \item We provide spectral leakage evidence showing that bidirectional allocation yields cleaner two-dimensional spectral structure after range compression.
\end{itemize}

\section{Method}

\subsection{Fixed-Budget Range-Azimuth Upsampling}

Let $s(\tau,\eta)$ denote the raw SAR echo, where $\tau$ and $\eta$ are the range-time and azimuth-time dimensions, respectively. In one-bit SAR imaging, the sampled echo is quantized by retaining only the sign information, which can be written as
\begin{equation}
    y = \operatorname{sign}(s + u),
\end{equation}
where $u$ denotes the quantization threshold. This operation greatly reduces the sampling precision but also discards amplitude information and introduces strong quantization distortion.

This paper studies how an upsampling budget should be allocated before one-bit quantization. Let $R$ be the range upsampling factor and $A$ be the azimuth upsampling factor. The total upsampling budget is defined as
\begin{equation}
    Q = R \times A .
\end{equation}
For a fixed $Q$, different allocation strategies can be constructed. A range-only strategy uses $R=Q$ and $A=1$, while an azimuth-only strategy uses $R=1$ and $A=Q$. A bidirectional strategy uses $R>1$ and $A>1$ under the same total budget. The no-upsampling case $R=1,A=1$ is used as the baseline.

The proposed bidirectional range-azimuth upsampling (BRAU) strategy does not change the SAR imaging processor itself. Instead, it changes the dimensional allocation of the sampling budget before one-bit quantization. After upsampling and quantization, all groups are processed by the same SAR imaging chain, including range compression, range cell migration correction, and final image formation.

\subsection{Threshold Settings}

We evaluate three threshold settings. Zero threshold (ZT) directly applies sign quantization with $u=0$. Non-random constant threshold (NCT) uses a deterministic constant threshold level. Random threshold (RT) uses a random threshold with an amplitude controlled by the scaling factor $A_s$. For RT, we also compare two implementation variants, SplitRT and FullRT, to test whether the observed allocation effect depends on the specific random-threshold construction.

Unless otherwise specified, the main fixed-budget experiments use SplitRT with $A_s=0.6$. This setting is used only to provide a consistent main comparison; the threshold robustness experiments verify that the bidirectional advantage is preserved under ZT, NCT, and RT.

\subsection{Evaluation Protocol}

Experiments are conducted on seven SAR datasets, with 10 frames selected from each dataset, giving 70 paired test samples. For each allocation group, image quality is evaluated by peak signal-to-noise ratio (PSNR) and structural similarity (SSIM). For each budget $Q$, we compare the best bidirectional allocation with the best unidirectional allocation. Statistical significance is tested using the paired Wilcoxon signed-rank test on the paired sample results.

\section{Experiments and Results}

\subsection{Main Fixed-Budget Results}

The main experiment evaluates whether splitting a fixed upsampling budget across range and azimuth is more effective than concentrating the same budget in a single dimension. We test integer budgets $Q=4,6,8,9,10$. For each $Q$, all available integer factor pairs are grouped into no-upsampling, range-only, azimuth-only, and bidirectional allocations.

\begin{figure}[t]
\centerline{\includegraphics[width=\linewidth]{../Exp1_MainResult_Output/Exp1_MainResult_PSNR_curves.png}}
\caption{Main fixed-budget PSNR results under different range-azimuth allocation strategies.}
\label{fig:main_psnr}
\end{figure}

\begin{figure}[t]
\centerline{\includegraphics[width=\linewidth]{../Exp1_MainResult_Output/Exp1_MainResult_SSIM_curves.png}}
\caption{Main fixed-budget SSIM results under different range-azimuth allocation strategies.}
\label{fig:main_ssim}
\end{figure}

Table~\ref{tab:main_gain} compares the absolute reconstruction quality of the best bidirectional allocation and the best unidirectional allocation under each budget. The last two columns report the gain, computed as bidirectional minus unidirectional. Across all tested budgets, the best bidirectional group consistently outperforms the best unidirectional group in both PSNR and SSIM. The PSNR improvement ranges from 0.2856 dB to 0.5072 dB, and the SSIM improvement ranges from 0.0169 to 0.0221. All paired Wilcoxon tests are statistically significant.

\begin{table}[t]
\caption{Best bidirectional and unidirectional results under each fixed budget.}
\begin{center}
\scriptsize
\setlength{\tabcolsep}{2pt}
\begin{tabular}{c c c c c}
\hline
$Q$ & Best bi-dir. & Best uni-dir. & $\Delta$PSNR & $\Delta$SSIM \\
\hline
4  & R2A2 (26.0623/0.8026) & R4A1 (25.7767/0.7857)  & 0.2856 & 0.0169 \\
6  & R2A3 (26.5585/0.8246) & R1A6 (26.1415/0.8037)  & 0.4170 & 0.0208 \\
8  & R2A4 (26.8194/0.8363) & R8A1 (26.3618/0.8148)  & 0.4576 & 0.0215 \\
9  & R3A3 (26.9167/0.8400) & R9A1 (26.4482/0.8186)  & 0.4685 & 0.0214 \\
10 & R5A2 (26.9702/0.8427) & R10A1 (26.4630/0.8206) & 0.5072 & 0.0221 \\
\hline
\end{tabular}
\label{tab:main_gain}
\end{center}
\end{table}

Figs.~\ref{fig:main_psnr} and \ref{fig:main_ssim} show the PSNR and SSIM curves. The no-upsampling baseline gives the lowest reconstruction quality. Increasing the total budget improves all upsampled groups, but the bidirectional groups form a consistently higher curve than the range-only and azimuth-only baselines. This result supports the central claim that the dimensional allocation of the upsampling budget matters in one-bit SAR imaging.

\subsection{Non-Integer Factorization}

To verify that the bidirectional advantage is not an artifact of integer factorization, we further evaluate non-integer allocation groups. The tested budgets include $Q=3,4.5,5,7.5$. In all cases, the best bidirectional allocation remains better than the best unidirectional allocation. The PSNR gains are 0.2723 dB, 0.3437 dB, 0.3952 dB, and 0.4567 dB, respectively. The corresponding SSIM gains range from 0.0146 to 0.0221. These results show that BRAU is not limited to balanced or integer-only allocation patterns.

\subsection{Mechanism Analysis Through Spectral Leakage}

We further analyze why bidirectional allocation is more effective under the same total budget. The working hypothesis is that SAR echoes have a two-dimensional spectral structure, while one-bit quantization introduces distortion across both range and azimuth. A unidirectional allocation increases sampling density in only one dimension, whereas bidirectional allocation distributes the budget over both dimensions and better suppresses spectral leakage after SAR processing.

For this analysis, we use $Q=4$ and compare four representative allocations: no upsampling (R1A1), range-only upsampling (R4A1), azimuth-only upsampling (R1A4), and bidirectional upsampling (R2A2). Fig.~\ref{fig:mechanism} visualizes the range-compressed spectra. The figure preserves the different spectral extents induced by the four allocation strategies, showing the one-dimensional expansion of R4A1 and R1A4 and the two-dimensional expansion of R2A2. Compared with the unidirectional groups, R2A2 produces a cleaner spectral concentration with lower off-support and directional leakage.

\begin{figure*}[t]
\centerline{\includegraphics[width=0.95\textwidth]{../Exp2_Mechanism_Output/Exp2_Node2_RC_Spectra.png}}
\caption{Range-compressed spectra for $Q=4$ under different allocation strategies. R4A1 and R1A4 expand the spectrum along one dimension, whereas R2A2 distributes the expansion over both range and azimuth and yields cleaner two-dimensional spectral concentration.}
\label{fig:mechanism}
\end{figure*}

\begin{table}[t]
\caption{Spectral leakage metrics after SAR processing for $Q=4$. Lower values indicate cleaner spectral concentration.}
\begin{center}
\scriptsize
\begin{tabular}{c c c c c}
\hline
Node & Allocation & Off-support & Range leak. & Azimuth leak. \\
\hline
RC & R1A1 & 0.926008 & 0.471044 & 0.101785 \\
RC & R4A1 & 0.893954 & 0.368013 & 0.070476 \\
RC & R1A4 & 0.901018 & 0.394427 & 0.072251 \\
RC & R2A2 & \textbf{0.892857} & \textbf{0.363750} & \textbf{0.066680} \\
\hline
RCMC & R1A1 & 0.926008 & 0.471044 & 0.101785 \\
RCMC & R4A1 & 0.893954 & 0.368013 & 0.070476 \\
RCMC & R1A4 & 0.901018 & 0.394427 & 0.072251 \\
RCMC & R2A2 & \textbf{0.892857} & \textbf{0.363750} & \textbf{0.066680} \\
\hline
\end{tabular}
\label{tab:mechanism_metrics}
\end{center}
\end{table}

Table~\ref{tab:mechanism_metrics} quantifies the visual trend in Fig.~\ref{fig:mechanism}. After RC, R2A2 achieves the lowest off-support ratio, range leakage ratio, and azimuth leakage ratio among all upsampled groups. The same ordering is retained after RCMC. This result supports the interpretation that bidirectional allocation better matches the two-dimensional range-azimuth structure of SAR data after coherent processing. The mechanism evidence is therefore consistent with the final reconstruction metrics: the bidirectional group is not only better in PSNR and SSIM, but also produces cleaner intermediate spectra.

\subsection{Robustness to Threshold Design}

Finally, we test whether the bidirectional advantage depends on a specific threshold construction. First, SplitRT and FullRT are compared over $A_s \in [0,1]$ for representative budgets. Their PSNR and SSIM curves nearly overlap, with only small differences across the tested range. This shows that the main effect is not tied to the internal implementation of random threshold generation.

\begin{figure}[t]
\centerline{\includegraphics[width=\linewidth]{../Exp3B_ZT_NCT_RT_Output/Exp3B_GainBars.png}}
\caption{Bidirectional gains under ZT, NCT, and RT threshold settings.}
\label{fig:threshold}
\end{figure}

Second, we compare ZT, NCT, and RT. For $Q=4,6,9$, the bidirectional group remains better than the best unidirectional group under all three threshold settings. Under ZT, the PSNR gains are 0.3612 dB, 0.4162 dB, and 0.4318 dB. Under NCT, the gains increase to 0.4308 dB, 0.4896 dB, and 0.5341 dB. Under RT, the gains remain positive at 0.2018 dB, 0.3306 dB, and 0.3611 dB. Therefore, RT improves the absolute reconstruction quality, but it is not the source of the bidirectional allocation effect.

\section{Conclusion}

This paper investigated fixed-budget range-azimuth upsampling allocation for one-bit SAR imaging. Instead of asking only how much to upsample or how to design the quantization threshold, we studied how a fixed total budget $Q=R\times A$ should be distributed between the range and azimuth dimensions. The experimental results show that bidirectional allocation consistently outperforms the best unidirectional allocation across all tested integer budgets. The same advantage is also observed under non-integer factors, indicating that the effect is not caused by a special integer factorization pattern.

The mechanism analysis provides further support for the allocation effect. After range compression and range cell migration correction, the bidirectional group produces lower off-support spectral energy and lower directional leakage than range-only and azimuth-only allocations. This suggests that splitting the budget across both dimensions better matches the two-dimensional spectral structure of SAR echoes under one-bit quantization. Threshold experiments also show that the advantage remains under ZT, NCT, and RT settings, so the observed effect is not tied to a specific threshold construction.

These findings support BRAU as a simple allocation strategy for improving one-bit SAR imaging under a fixed sampling budget. Future work will extend the analysis to broader SAR scenes and develop a more explicit theoretical model for two-dimensional quantization noise shaping in range-azimuth processing.

\begin{thebibliography}{00}
\bibitem{cumming2005sar}
I. G. Cumming and F. H. Wong, \textit{Digital Processing of Synthetic Aperture Radar Data: Algorithms and Implementation}. Norwood, MA, USA: Artech House, 2005.

\bibitem{curlander1991sar}
J. C. Curlander and R. N. McDonough, \textit{Synthetic Aperture Radar: Systems and Signal Processing}. New York, NY, USA: Wiley, 1991.

\bibitem{franceschetti1991signum}
G. Franceschetti, V. Pascazio, and G. Schirinzi, ``Processing of signum coded SAR signal: Theory and experiments,'' \textit{IEE Proceedings F Radar and Signal Processing}, vol. 138, no. 3, pp. 192--198, 1991, doi: 10.1049/ip-f-2.1991.0025.

\bibitem{zhao2019onebit}
B. Zhao, L. Huang, and W. Bao, ``One-bit SAR imaging based on single-frequency thresholds,'' \textit{IEEE Transactions on Geoscience and Remote Sensing}, vol. 57, no. 9, pp. 7017--7032, 2019, doi: 10.1109/TGRS.2019.2910284.

\bibitem{demir2018onebit}
M. Demir and E. Ercelebi, ``One-bit compressive sensing with time-varying thresholds in synthetic aperture radar imaging,'' \textit{IET Radar, Sonar \& Navigation}, vol. 12, no. 12, pp. 1517--1526, 2018, doi: 10.1049/iet-rsn.2018.5044.

\bibitem{zhao2021lowprecision}
B. Zhao, L. Huang, and B. Jin, ``Strategy for SAR imaging quality improvement with low-precision sampled data,'' \textit{IEEE Transactions on Geoscience and Remote Sensing}, vol. 59, no. 4, pp. 3150--3160, 2021, doi: 10.1109/TGRS.2020.3014300.

\bibitem{nie2025fixed}
G. Nie, B. Zhao, Q. Liu, L. Huang, and G. Liao, ``One-bit synthetic aperture radar imaging based on fixed-threshold with slow-time fluctuations,'' \textit{IEEE Transactions on Geoscience and Remote Sensing}, vol. 63, pp. 1--15, 2025, doi: 10.1109/TGRS.2024.3519757.

\bibitem{si2023dequantization}
C. Si, B. Zhao, L. Huang, and S. Liu, ``A convolutional de-quantization network for harmonics suppression in one-bit SAR imaging,'' \textit{IEEE Transactions on Geoscience and Remote Sensing}, vol. 61, pp. 1--16, 2023, doi: 10.1109/TGRS.2023.3330530.
\end{thebibliography}

\end{document}
